\documentclass[10pt,a4paper]{article}

% --- PACKAGES ---
\usepackage[utf8]{inputenc}
\usepackage[T1]{fontenc}
\usepackage[scaled=0.92]{helvet}
\usepackage{mathpazo}
\usepackage[margin=1in]{geometry}
\usepackage{amsmath,amssymb}
\usepackage{booktabs}
\usepackage{tabularx}
\usepackage{multicol}
\usepackage{enumitem}
\usepackage{titlesec}
\usepackage{fancyhdr}
\usepackage{xcolor}
\usepackage{listings}
\usepackage{tcolorbox}
\usepackage[hidelinks]{hyperref}

% --- SETUP ---
\setlength{\parindent}{0pt}
\setlength{\parskip}{0.4em}
\setlist[itemize]{leftmargin=*, noitemsep, topsep=0pt, parsep=0pt}
\setlist[enumerate]{leftmargin=*, noitemsep, topsep=0pt, parsep=0pt}

% --- COLORS & STYLES ---
\definecolor{codeblue}{rgb}{0.1,0.2,0.6}
\definecolor{codegray}{rgb}{0.4,0.4,0.4}
\definecolor{codebg}{rgb}{0.96,0.96,0.96}

\titleformat{\section}{\normalfont\large\bfseries\sffamily\color{codeblue}}{}{0em}{}[\titlerule]
\titlespacing*{\section}{0pt}{1.5ex}{1ex}

\titleformat{\subsection}{\normalfont\small\bfseries\sffamily}{}{0em}{}
\titlespacing*{\subsection}{0pt}{1ex}{0.5ex}

\lstset{
    basicstyle=\ttfamily\footnotesize,
    keywordstyle=\color{codeblue}\bfseries,
    backgroundcolor=\color{codebg},
    frame=single,
    rulecolor=\color{codegray!30},
    breaklines=true,
    aboveskip=2pt,
    belowskip=2pt
}

% Header & Footer
\pagestyle{fancy}
\fancyhf{}
\fancyhead[L]{\small\sffamily\bfseries Cheat Sheet Fondasi Basis Data (P1-P7)}
\fancyhead[R]{\small\sffamily Muhammad Riduwan Khafidi}
\renewcommand{\headrulewidth}{0.4pt}

% Custom box for Tips
\newtcolorbox{tipbox}{
    colback=gray!5, colframe=gray!40,
    boxrule=0.5pt, arc=2pt,
    left=4pt, right=4pt, top=4pt, bottom=4pt,
    fontupper=\small\sffamily
}

\begin{document}

\begin{multicols*}{2}

	% ==========================================
	% BAGIAN 1: KONSEP DASAR & ERD
	% ==========================================
	\section{Konsep Dasar & ERD (P1-P3)}

	\textbf{Data vs Database:} Data = fakta mentah. Database = kumpulan data terorganisir yang saling berelasi. DBMS = software pengelola (MySQL, Oracle, dll).

	\textbf{Komponen ERD:}
	\begin{itemize}
		\item \textbf{Entitas (Kotak):} Objek nyata (Pegawai, Barang).
		\item \textbf{Atribut (Elips):} Properti entitas.
		      \begin{itemize}
			      \item \textit{Primary Key:} Digarisbawahi.
			      \item \textit{Multivalued:} Elips ganda (No. Telp).
			      \item \textit{Derivatif:} Elips putus-putus (Umur dari Tgl Lahir).
		      \end{itemize}
		\item \textbf{Relasi (Belah Ketupat):} Hubungan antar entitas.
		\item \textbf{Weak Entity (Kotak Ganda):} Entitas tanpa PK, butuh entitas kuat.
	\end{itemize}

	\textbf{Kardinalitas Relasi:}
	\begin{itemize}
		\item \textbf{1:1} (Satu Dept punya satu Manajer).
		\item \textbf{1:N} (Satu Dept punya banyak Pegawai).
		\item \textbf{N:M} (Banyak Mahasiswa ambil banyak Matkul).
	\end{itemize}

	\textbf{Participation Constraints:}
	\begin{itemize}
		\item \textit{Total (Garis Ganda):} Wajib berpartisipasi (Pinjaman wajib punya nasabah).
		\item \textit{Partial (Garis Tunggal):} Boleh tidak berpartisipasi.
	\end{itemize}

	\textbf{EERD (Enhanced ERD):}
	\begin{itemize}
		\item \textbf{Specialization/Generalization:} Sub-grouping (Dosen $\rightarrow$ Dosen Tetap / Tidak Tetap). Simbol segitiga "IS A".
		\item \textbf{Aggregation:} Relasi diperlakukan sebagai entitas baru.
	\end{itemize}

	% ==========================================
	% BAGIAN 2: MODEL RELASIONAL & KONVERSI
	% ==========================================
	\section{Model Relasional (P4-P5)}
	\textbf{Istilah:} Relasi = Tabel, Tuple = Baris/Record, Atribut = Kolom/Field. Degree = jumlah kolom. Cardinality = jumlah baris.

	\textbf{Jenis Key:}
	\begin{itemize}
		\item \textbf{Super Key:} Kumpulan atribut pembeda.
		\item \textbf{Candidate Key:} Super key paling minimal.
		\item \textbf{Primary Key (PK):} Candidate key yang dipilih.
		\item \textbf{Foreign Key (FK):} Atribut yang merujuk PK tabel lain.
	\end{itemize}

	\textbf{Konversi ERD ke Tabel:}
	\begin{enumerate}
		\item \textbf{Entitas Kuat:} Jadi Tabel. Atribut jadi kolom, PK jadi Primary Key.
		\item \textbf{Entitas Lemah:} Jadi Tabel. Tambahkan PK entitas kuat sebagai FK (wajib \texttt{ON DELETE CASCADE}).
		\item \textbf{Relasi 1:1:} Tambahkan FK di salah satu tabel.
		\item \textbf{Relasi 1:N:} Taruh FK di sisi "Many" (Banyak).
		\item \textbf{Relasi N:M:} Buat \textbf{Tabel Baru} berisi PK dari kedua entitas (sebagai Composite PK).
		\item \textbf{Multivalued:} Buat tabel baru (PK entitas asli + atribut multivalued).
	\end{enumerate}

	% ==========================================
	% BAGIAN 3: DDL (P6)
	% ==========================================
	\section{DDL - Membuat Struktur (P6)}
	\textbf{Tipe Data Penting:}
	\begin{itemize}
		\item \texttt{CHAR(n)}: Panjang tetap (boros memori, tapi cepat).
		\item \texttt{VARCHAR(n)}: Panjang dinamis (hemat memori).
		\item \texttt{INT}, \texttt{NUMERIC(p,d)}, \texttt{DATE}, \texttt{NULL}.
	\end{itemize}

	\textbf{Create \& Alter Table:}
	\begin{lstlisting}[language=SQL]
CREATE TABLE branch (
  branch_name VARCHAR(15) NOT NULL,
  assets INT,
  PRIMARY KEY (branch_name),
  CHECK (assets >= 0)
);

ALTER TABLE borrower
  ADD b_date DATE; -- Tambah kolom
ALTER TABLE borrower
  CHANGE b_date birth_date DATE; -- Ganti nama
ALTER TABLE borrower
  DROP birth_date; -- Hapus kolom
\end{lstlisting}

	\textbf{Constraints Penting:}
	\begin{itemize}
		\item \texttt{NOT NULL}, \texttt{UNIQUE}, \texttt{CHECK (kondisi)}.
		\item \texttt{FOREIGN KEY ... REFERENCES ... ON DELETE CASCADE}: Jika data induk dihapus, data anak ikut terhapus.
	\end{itemize}

	\textbf{Index (Kecepatan Pencarian):}
	\begin{lstlisting}[language=SQL]
CREATE INDEX branch_idx ON Account (branch_name);
\end{lstlisting}
	\textbf{CRUD (DML Dasar):}
	\begin{lstlisting}[language=SQL]
INSERT INTO account VALUES ('A-1', 'Batam', 2000);
DELETE FROM account WHERE branch_name = 'Batam';
UPDATE account SET balance = balance * 1.05 WHERE balance > 1000;
\end{lstlisting}

	% ==========================================
	% BAGIAN 4: DML (P7)
	% ==========================================
	\section{DML - Query \& Manipulasi (P7)}
	\textbf{Struktur Dasar:}
	\begin{lstlisting}[language=SQL]
SELECT A1, A2, ... FROM R1, R2 WHERE P;
\end{lstlisting}

	\textbf{Operator WHERE:}
	\begin{itemize}
		\item Logika: \texttt{AND, OR, NOT}.
		\item Perbandingan: \texttt{=, >, >=, <, <=, <>}.
		\item String: \texttt{LIKE} (\% = banyak karakter, \_ = 1 karakter).
		\item Cek Kosong: \texttt{IS NULL} / \texttt{IS NOT NULL}.
		\item Rentang: \texttt{IN ('102', '104')}, \texttt{BETWEEN 10 AND 20}.
	\end{itemize}

	\textbf{Fungsi Agregat (Hasil 1 Baris):}
	\texttt{AVG, MIN, MAX, SUM, COUNT}.
	\begin{lstlisting}[language=SQL]
SELECT COUNT(*) FROM Customer;
SELECT MIN(Balance), MAX(Balance) FROM Account;
\end{lstlisting}

	\textbf{Join 2 Tabel (FROM Extension):}
	Menggabungkan 2 tabel dengan WHERE (menghindari Cartesian Product):
	\begin{lstlisting}[language=SQL]
SELECT A.Owner, A.Balance
FROM Account A, Deposit D
WHERE D.Account = A.Number AND D.Amount > 1000;
\end{lstlisting}

	\textbf{CASE Statement (Percabangan):}
	\begin{lstlisting}[language=SQL]
UPDATE account SET balance =
  CASE
    WHEN balance <= 1000 THEN balance * 1.05
    ELSE balance * 1.06
  END;
\end{lstlisting}

	% ==========================================
	% BAGIAN 5: TIPS & TRIK (CONTEKAN)
	% ==========================================
	\section*{💡 Tips \& Trik Ujian / Praktikum}

	\begin{tipbox}
		\textbf{1. Hati-hati dengan NULL!}\\
		Jangan pernah menulis \texttt{WHERE Address = NULL}. SQL tidak akan pernah mengembalikan hasil. Gunakan \texttt{WHERE Address IS NULL}.
	\end{tipbox}

	\begin{tipbox}
		\textbf{2. Kapan pakai CHAR vs VARCHAR?}\\
		Gunakan \texttt{CHAR} jika panjang datanya pasti (misal: Kode Negara 'ID', 'US'). Gunakan \texttt{VARCHAR} jika panjangnya bervariasi (misal: Nama, Alamat).
	\end{tipbox}

	\begin{tipbox}
		\textbf{3. Bahaya Cartesian Product!}\\
		Jika menulis \texttt{SELECT * FROM A, B} tanpa klausa \texttt{WHERE} yang menghubungkan PK dan FK, hasilnya adalah perkalian semua baris A dan B (Ribuan data sampah).
	\end{tipbox}

	\begin{tipbox}
		\textbf{4. Kapan Index digunakan?}\\
		Gunakan \texttt{INDEX} HANYA pada kolom yang sering dicari (\texttt{WHERE kolom = ...}). \textbf{JANGAN} gunakan pada tabel kecil atau tabel yang sering di-\textit{update}, karena setiap update akan memperlambat kinerja karena index harus diperbarui.
	\end{tipbox}

	\begin{tipbox}
		\textbf{5. Ingat Aturan Konversi N:M!}\\
		Relasi Many-to-Many (N:M) \textbf{TIDAK BISA} langsung dimasukkan sebagai Foreign Key di salah satu tabel. Wajib dibuatkan \textbf{Tabel Baru} (Tabel Penghubung/Bridge Table).
	\end{tipbox}

	\begin{tipbox}
		\textbf{6. Perbedaan DELETE vs DROP:}\\
		\texttt{DELETE FROM table} = Menghapus data (baris). Tabel masih ada.\\
		\texttt{DROP TABLE table} = Menghapus struktur tabel secara permanen beserta datanya.
	\end{tipbox}

\end{multicols*}
\end{document}
